\begin{table}[htbp]\centering
\caption{Approximate minimum detectable effects for selected specifications}
\label{tab:mde}
\small
\resizebox{\textwidth}{!}{%
\begin{tabular}{lcccccc}
\toprule
 & estimate & SE & MDE (80\% power) & MDE / SD of outcome & units & N \\
\midrule
H1a: labour share on log K/L & -0.094 & 0.051 & 0.142 & 1.84 & share units (multiply by 100 for pp) & 555 \\
H1b: mobilization on log(sK/sL), low repression & 0.446 & 0.225 & 0.629 & 0.70 & mobilization units & 292 \\
H2: physical repression on log K/L & -0.027 & 0.079 & 0.221 & 1.11 & index units (SD 0.21) & 569 \\
H2: physical repression on K/L x anti-system threat & -0.274 & 0.307 & 0.861 & 4.34 & index units & 569 \\
H2: redistribution on lagged repression & -2.184 & 4.449 & 12.457 & 3.83 & Gini points & 553 \\
H3: log robots per worker on log K/L & 0.604 & 0.262 & 0.734 & 0.94 & log points & 461 \\
\bottomrule
\end{tabular}}
\begin{tablenotes}\footnotesize Normal approximation: MDE = 2.8 x country-clustered SE for 80 percent power at a two-sided 5 percent level. Outcomes retain original units; continuous regressors are standardized. With 18 clusters this is a sensitivity diagnostic, not a finite-sample power guarantee or an equivalence test. MDE / SD divides by the outcome SD. For interactions the coefficient is on the standardized product, not a marginal effect. Specifications differ from the multiple-testing family.\end{tablenotes}
\end{table}